% fgib_bound.tex —— FGIB 逆正交变换界与 CEB 界的紧致性比较推导
% 记号沿用 Fischer (2020)、Kolchinsky et al. (2017, NIB)、Kolchinsky & Tracey (2019, Caveats)
% 与 fgib_degeneration.tex；用法同 fgib_linear.tex
% 核心结论（两个方向，各带条件）：
% (1) 变分值意义：同一编码器 q 下，CEB 可学习反向编码器把 gap 消到矩匹配最小，CEB 的界 ≤ FGIB 的界，
%     差额恰为锚点 gap —— 对齐族上即 ½(s-a)² + (z/2)(σ²-1-ln σ²)，与 fgib_degeneration.tex 的扫参项同一表达式；
% (2) 作用对象意义：经逆线性变换 A⁻¹（mu_head 无约束，正交性在锚点框架与平衡态），FGIB 的界作用在"推理表示 h + 已知旁路噪声"上；
%     CEB 的界作用在训练期噪声码 Z 上，由 DPI 严格低于推理表示 μ(X) 的残差信息——作为推理表示残差的上界，
%     FGIB 的界更紧（CEB 的界根本不是该对象的有效上界），且 FGIB 的缺口可由仅旁路参数 anchor_var→0 消去。
% (3) 正交因子条件数：锚点参考侧（QR 正交框架 + 单位协方差，固定 buffer）条件数恒为 1；
%     编码器侧度量 A^TA 可学习、全局条件数无构造保证，但对齐平衡族上限制在类均值子空间为 s²I
%     （有效方向条件数 1）、且 A 只有指向正交框架的对齐梯度、无退化驱动力（正交参数化 A 为可选硬保证）；
%     CEB 的 MLP 先验两侧皆可训练：协方差条件数可达 e^{20}、均值可秩亏、方差竞赛把界推向 e^{10} 权重退化区
%     （prop:conditioning；两种"更紧"的语义见 rem:senses）。

\newtheorem{lemma}{Lemma}
\newtheorem{proposition}{Proposition}
\newtheorem{remark}{Remark}

\textbf{Notation.}
$X$ input, $Y$ label, $K$ classes.
Both models share the encoder posterior $q(z|x)$; write
$q(z|y) := \mathbb{E}_{p(x|y)}[q(z|x)]$ for the class-conditional marginal.
CEB's backward encoder is a learnable per-class (diagonal) Gaussian
$b_\theta(z|y) = \mathcal{N}(m_\theta(y), S_\theta(y))$ (linear \texttt{prior\_net} on
one-hot $y$); FGIB's prior is the fixed anchor $r(z|y) = \mathcal{N}(a\,v_y, I)$ with
orthonormal $v_y$ and anchor variance $1$.
The variational objects (Fischer 2020, Eqs.~18--20; Alemi et al.\ 2017, Eq.~14;
Kolchinsky et al.\ 2017, Eqs.~10 and 16):
\begin{align*}
R &= \mathbb{E}\big[\operatorname{KL}(q(z|x)\,\|\,m(z))\big] \ge I(X;Z)
&&\text{(VIB rate, marginal prior $m$)} \\
\mathrm{ReX} &= \mathbb{E}\big[\operatorname{KL}(q(z|x)\,\|\,b_\theta(z|y))\big]
\ge I(X;Z|Y)
&&\text{(CEB residual)} \\
\widehat I &= -\frac{1}{N}\sum_i \ln\Big[\frac{1}{N}\sum_j
e^{-\|f_i-f_j\|^2/(2\sigma^2)}\Big] \ge I(X;M)
&&\text{(NIB, non-parametric; Eq.~16, the homoscedastic case of Eq.~10)} \\
\mathrm{ReX}_F &= \mathbb{E}\big[\operatorname{KL}(q(z|x)\,\|\,r(z|y))\big]
\ge I(X;Z|Y)
&&\text{(FGIB, fixed anchor)}
\end{align*}
Fischer's comparison of the first two: at the respective prior optima
($m = q(z)$, $b = q(z|y)$), $\mathrm{ReX} = R - I(Y;Z)$ --- the conditional prior removes
exactly the retained label information, so CEB's bound is tighter than VIB's by
$I(Y;Z)$. We ask the same question one level down: where does FGIB's fixed-anchor bound
sit relative to CEB's learnable-backward bound?

\begin{lemma}[Gap identity]\label{lem:gap}
For any conditional $r(z|y)$,
\begin{equation}
\mathbb{E}\big[\operatorname{KL}(q(z|x)\,\|\,r(z|y))\big]
= I(X;Z|Y) + \mathbb{E}_y\big[\operatorname{KL}(q(z|y)\,\|\,r(z|y))\big] .
\label{eq:gap}
\end{equation}
\end{lemma}

\begin{proof}
$\mathbb{E}[\operatorname{KL}(q(z|x)\|r(z|y))] = \mathbb{E}[\ln q(z|x) - \ln r(z|y)]$
and $I(X;Z|Y) = \mathbb{E}[\ln q(z|x) - \ln q(z|y)]$ (Fischer 2020, Eqs.~9--10), where
both expectations are over $p(x,y)q(z|x)$; subtracting,
$\mathbb{E}[\ln q(z|y) - \ln r(z|y)] = \mathbb{E}_y[\operatorname{KL}(q(z|y)\|r(z|y))]$
since the argument depends on $z, y$ only. \hfill\qed
\end{proof}

\begin{lemma}[Moment matching]\label{lem:mm}
Within the per-class Gaussian (diagonal) family, the gap of Lemma~\ref{lem:gap} is
minimized by the moment-matched Gaussian $b^*_y = \mathcal{N}(m_y, S_y)$ with
$m_y = \mathbb{E}\,q(z|y)$, $S_y = \operatorname{diag}\operatorname{Var} q(z|y)$.
\end{lemma}

\begin{proof}
$\operatorname{KL}(p\|q)$ as a function of $q$ is minimized by matching the first two
moments of $p$ (standard; for Gaussians the minimizer is unique). \hfill\qed
\end{proof}

\begin{proposition}[Variational ordering: CEB's bound is tighter in value]\label{prop:ordering}
For every fixed encoder $q(z|x)$,
$\mathrm{ReX}(\theta^*) \le \mathrm{ReX}_F$ where $\theta^*$ is CEB's optimal backward
encoder, and the excess is the \emph{anchor gap}
\begin{equation}
\Delta := \mathrm{ReX}_F - \mathrm{ReX}(\theta^*)
= \mathbb{E}_y\big[\operatorname{KL}(q(z|y)\,\|\,N(a v_y, I))\big]
- \mathbb{E}_y\big[\operatorname{KL}(q(z|y)\,\|\,N(m_y, S_y))\big] \ge 0 ,
\label{eq:anchor-gap}
\end{equation}
with equality iff the anchors equal the moment-matched marginals.
On the aligned family of fgib\_degeneration.tex ($q(z|y) = \mathcal{N}(s v_y, \sigma^2 I)$, so
$I(X;Z|Y) = 0$ and $b^* = q(z|y)$ exactly),
\begin{equation}
\Delta = \tfrac{1}{2}(s-a)^2 + \tfrac{z}{2}\big(\sigma^2 - 1 - \ln\sigma^2\big) ,
\label{eq:aligned-gap}
\end{equation}
i.e.\ $\mathrm{ReX}_F = \Delta$ and $\mathrm{ReX}(\theta^*) = 0$: FGIB's bound value
\emph{equals} the sweep-driving compression term of fgib\_degeneration.tex
(its sweep proposition), while CEB's bound collapses to the (zero) residual.
\end{proposition}

\begin{proof}
The anchor $N(a v_y, I)$ is a member of CEB's per-class Gaussian family, so the minimum
over the family is no larger than the value at the anchor: the ordering and
Eq.~\eqref{eq:anchor-gap} follow from Lemmas~\ref{lem:gap} and~\ref{lem:mm}.
On the aligned family the marginal is itself Gaussian, so $b^* = q(z|y)$ gives gap $0$,
$I(X;Z|Y) = 0$ (the posterior is class-constant), and
$\operatorname{KL}(\mathcal{N}(s v_y, \sigma^2 I)\|\mathcal{N}(a v_y, I))
= \tfrac{1}{2}(s-a)^2 + \tfrac{z}{2}(\sigma^2-1-\ln\sigma^2)$. \hfill\qed
\end{proof}

\begin{remark}[Tightness--force trade-off]\label{rem:force}
Proposition~\ref{prop:ordering} is the price of the fixed anchor, and it is paid for a
reason: the gap is the only quantity the encoder can reduce. CEB's gap is a function of
the \emph{backward encoder alone} --- $b_\theta \to q(z|y)$ eliminates it with the encoder
held fixed, so CEB's bound can become tight while the representation stays uncompressed:
its tightness certifies nothing about compression, and the gap is an unobservable
optimization residual that tracks the forward encoder. FGIB's gap can shrink only if the
\emph{encoder} moves toward the anchors (the anchor is fixed), and on the aligned family
the bound value \emph{is} the compression certificate (Eq.~\eqref{eq:aligned-gap}): bound
tightness is monotonically equivalent to compression. This is the same fixed-anchor
mechanism that keeps the sweep force persistent (fgib\_degeneration.tex, its sweep proposition).
\end{remark}

\begin{proposition}[The inverse linear transform: the object of the bound]\label{prop:inverse}
Let the side head $A$ (FGIB's introduced linear layer, modeled without its bias by the
convention of fgib\_degeneration.tex) be \emph{invertible} ($z = d$, unconstrained), and
let the trained posterior be isotropic, $\Sigma(x) = \sigma_p^2 I$ (the KL variance term
is minimized at $\sigma^2 = \sigma_p^2$, the anchor variance, so this holds at the
variance optimum). Then $\tilde z := A^{-1} z$ gives
$q'(\tilde z|x) = \mathcal{N}(h(x), \sigma_p^2 (A^{\top}A)^{-1})$ --- the posterior is
centered at the \emph{inference representation} $h$ with side noise
$\sigma_p^2 (A^{\top}A)^{-1}$ --- and the transformed anchors
$r'(\tilde z|y) = \mathcal{N}(a\, A^{-1} v_y, (A^{\top}A)^{-1})$ form an orthonormal
frame at scale $a$ in the induced metric $A^{\top}A$
($\langle A^{-1}v_i, A^{-1}v_j\rangle_{A^{\top}A} = \delta_{ij}$: the fixed frame
$\{v_y\}$ read in canonical coordinates). By KL invariance under invertible linear maps
(fgib\_linear.tex, Lemma~1),
\begin{equation}
\mathrm{ReX}_F
= I(X;\, h + \sigma_p\,(A^{\top}A)^{-1/2}\varepsilon\,|Y)
+ \mathbb{E}_y\big[\operatorname{KL}\big(q'(\tilde z|y)\,\|\,r'(\tilde z|y)\big)\big],
\qquad \varepsilon \sim \mathcal{N}(0, I) .
\label{eq:object}
\end{equation}
The bound's object is the inference representation perturbed by \emph{known,
side-path-only} noise (its covariance is a model parameter), and the anchor geometry is
preserved exactly under the transform, in the induced metric.
\end{proposition}

\begin{proof}
$A^{-1}A = I$; $A^{-1}\Sigma A^{-\top} = \sigma_p^2 A^{-1}A^{-\top}
= \sigma_p^2 (A^{\top}A)^{-1}$ for $\Sigma = \sigma_p^2 I$; the prior covariance maps to
$A^{-1} I A^{-\top} = (A^{\top}A)^{-1}$; and
$\langle A^{-1}v_i, A^{-1}v_j\rangle_{A^{\top}A}
= v_i^{\top} A^{-\top} A^{\top} A\, A^{-1} v_j = \delta_{ij}$.
Equation~\eqref{eq:object} is Lemma~\ref{lem:gap} applied to the transformed pair. \hfill\qed
\end{proof}

\begin{proposition}[Tightness on the inference object: the sense in which FGIB's bound is tighter]\label{prop:object}
(a)~\emph{CEB}. The classifier consumes the sampled code $Z$ at training and the mean
$\mu(X)$ at evaluation (Fischer 2020, Appendix~A.1). At the variance-optimal posterior,
$\sigma(x) \equiv \sqrt{S_\theta(y)}$ (the backward encoder's per-class standard
deviation) is class-constant, so $Z = \mu(X) + \sqrt{S_\theta(Y)}\,\varepsilon$ is a
noisy function of $\mu(X)$ given $Y$; by the DPI,
\begin{equation}
I(X; Z | Y) \;\le\; I(X;\, \mu(X) | Y) ,
\label{eq:dpi}
\end{equation}
strictly so whenever the noise destroys within-class detail. Hence CEB's bound
$\mathrm{ReX} \ge I(X;Z|Y)$ is \emph{not} an upper bound on the residual information of
the representation actually used at inference: its object sits strictly below the
inference residual, and the deficit $I(X;\mu(X)|Y) - I(X;Z|Y)$ --- the within-class
information destroyed by the training noise --- is positive, a trained quantity
($S_\theta$ is a model output), and unobservable during training. The CE term presses the
noise toward $0$ (by Jensen,
$\mathbb{E}_z[\mathrm{CE}(c(z),y)] \ge \mathrm{CE}(c(\mu),y)$): this shrinks the
deficit, but it drives the pair into the noise-free regime where the certificate
degenerates (fgib\_linear.tex, noise-free-limit remark), so the deficit can never be read
off the certificate.
(b)~\emph{FGIB}. The inference path is $h$ at both training and evaluation; the side
noise of Proposition~\ref{prop:inverse} is not part of any inference path. The bound's
deficit $I(X;H|Y) - I(X;\,h+\sigma_p(A^{\top}A)^{-1/2}\varepsilon\,|Y)$ is the
information destroyed by a known perturbation, and it is removable by
\emph{side-path-only} parameters:
as the anchor variance $\sigma_p^2 \to 0$ (point anchors), the variance-optimal posterior
collapses ($\sigma^2 = \sigma_p^2 \to 0$) and the bound's object converges to
$I(X;H|Y)$ itself, without changing the classifier or the inference behavior.
CEB cannot take this limit: its noise is in the training path, and removing it changes
the model (the classifier is trained on the noisy codes).
(c)~\emph{Conclusion}. In the pure variational sense (same encoder, same Gaussian
family), CEB's bound is tighter (Proposition~\ref{prop:ordering}). As an upper bound on
the residual of the representation used at inference --- the object the regularizer is
meant to control --- FGIB's bound, taken through the inverse linear transform, is the
tighter (in the point-anchor limit, the only valid) one: its object is the inference
representation up to a known, removable perturbation, while CEB's object is the training
code, which lies strictly below the inference residual by an unobservable, model-driven
deficit (Eq.~\eqref{eq:dpi}).
\end{proposition}

\begin{remark}[The bound is a training instrument, not an inference one]\label{rem:training-role}
No method computes its bound at inference: the bound's roles are \emph{training-time}
--- a regularizer (a differentiable surrogate for the intractable information term) and
a meter (Fischer 2020, p.~7: ``observing ReX tells us exactly how many bits we are from
the MNI point''). In the $Z \perp X \mid Y$ limit CEB's mean is also class-constant
($m(x) = \mathbb{E}[Z|x] = \mathbb{E}[Z|y]$), so the mismatch of
Prop.~\ref{prop:object}(a) is not a limit failure: it is a finite-$\beta$
\emph{under-reporting} --- the regularizer controls, and the meter reports, the training
code, which is strictly more compressed than the deployed representation $\mu(X)$, by a
positive, trained, unobservable deficit. FGIB's side-path noise is known and removable
($\sigma_p \to 0$), so its regularizer and meter act on the deployed object itself.
The one genuine inference-time use of a certificate value is as a detection score
(Fischer 2020, RG3, uses the marginal rate $R$): CEB's label-conditioned
$\mathrm{ReX}$ is not even computable at inference (the label is unknown --- Fischer
trains a separate marginal for this purpose), while FGIB's geometry gives a label-free
score, the distance to the nearest anchor
$\min_y \operatorname{KL}(q(z|x)\|\mathcal{N}(a v_y, I))$, with no extra network
(the minimum is an internal enumeration over the $K$ fixed anchor buffers --- the single
QR orthonormal frame --- and consults no ground-truth label, no more than softmax
consults labels when it enumerates the $K$ logits).
\end{remark}

\begin{proposition}[Orthogonal-factor conditioning: the reference is fixed and well-conditioned]\label{prop:conditioning}
FGIB's anchor frame $V = [v_1, \dots, v_K]$ is the orthonormal factor of a random matrix
(QR: $G = QR$, $V := Q$, $V^{\top}V = I_K$), and the anchor covariance is the identity.
Then, for every anchor scale $a \ge 0$ and throughout training:
\begin{enumerate}
\item[(i)] the \emph{reference} of the matching --- the prior's mean frame and
covariance --- is fixed (non-trainable buffers) and has condition number $1$: the
certificate is a distance to a compact, non-degenerate, class-separated (for $a > 0$)
target;
\item[(ii)] in side-mean coordinates the certificate's displacement metric is the
identity: the KL's displacement term is $\tfrac12\|\mu - a v_y\|^2$ with gradient
$\mu - a v_y$;
\item[(iii)] on the shared-encoder side the metric is $A^{\top}A$, where $A$ is a
learnable unconstrained map, so its \emph{global} condition number is not controlled by
construction. However, on the class-aligned equilibrium family ($A h_y = s v_y$, the
family of fgib\_degeneration.tex) the metric restricted to the class-mean subspace is
$s^2 I$ --- condition number $1$ on every direction that carries the matching force ---
and the only gradient $A$ receives is the alignment force
$(\mu - a v_y)\,h^{\top}$ toward the orthonormal target, which never points toward
degeneracy; an orthogonal parameterization of $A$ is the optional hard guarantee
(Remark~\ref{rem:orth}).
\end{enumerate}
CEB's backward encoder $b_\theta(z|y) = \mathcal{N}(m_\theta(y), S_\theta(y))$ is an
unconstrained linear map of the label: \emph{both} sides of its matching are trainable,
its covariance diagonal can lie anywhere in the clamp $[e^{-10}, e^{10}]$ (condition
number up to $e^{20}$), its means can be rank-deficient (class-merged codes, or the
collapse solution $m_\theta \equiv 0$), and the variance race drives $S_\theta$ to the
clamp, degenerating the certificate into a mean-matching force of weight $1/s = e^{10}$
(fgib\_ceb\_collapse.tex, Prop.~variance). The fixed QR reference is what keeps FGIB's
bound well-conditioned in the regime where CEB's bound degenerates.
\end{proposition}

\begin{proof}
$V^{\top}V = I_K$: all singular values of $V$ are $1$; the covariance is $I$ with all
eigenvalues $1$; both are frozen buffers, so (i) holds. The diagonal closed form of
$\operatorname{KL}(\mathcal{N}(\mu, \sigma^2 I)\,\|\,\mathcal{N}(a v_y, I))$ has
displacement term $\tfrac12\|\mu - a v_y\|^2$, giving (ii). For (iii):
$h_i^{\top} A^{\top} A\, h_j = \mu_i^{\top}\mu_j = s^2 v_i^{\top} v_j = s^2\delta_{ij}$
on the aligned family, and the KL gradient on $A$ is
$(\mu - a v_y)\,h^{\top}$, an alignment force toward the orthonormal frame that vanishes
only at perfect alignment. For CEB, the displacement term is
$\tfrac12(m - g)^{\top} S_\theta^{-1}(m - g)$: the metric $S_\theta^{-1}$ has condition
number up to $e^{20}$ under the logvar clamp, and the gradient weight $1/s$ at the
collapse equilibrium is fgib\_ceb\_collapse.tex, Prop.~variance. \hfill\qed
\end{proof}

\begin{remark}[Two senses of tightness]\label{rem:senses}
In the pure variational \emph{value} sense (same encoder, same Gaussian family), CEB's
bound is tighter: $\mathrm{ReX}(\theta^*) \le \mathrm{ReX}_F$
(Prop.~\ref{prop:ordering}), the excess being the anchor gap --- the price of the fixed
anchor. The statement ``FGIB's bound is tighter'' holds in the \emph{object} sense: as
an upper bound on the residual of the representation used at inference, CEB's optimized
bound is not a bound at all --- it can lie strictly below its target
(Prop.~\ref{prop:object}(a)) --- while FGIB's bound is valid, has a fixed
well-conditioned reference and a closed variance channel
(Prop.~\ref{prop:conditioning}), and becomes exact in the point-anchor limit
(Prop.~\ref{prop:object}(b)). Being below the target is not tightness; it is invalidity.
The learnability of the side head $A$ does not disturb this picture: $A$ is part of the
encoder, so the value ordering of Prop.~\ref{prop:ordering} holds for every fixed $q$,
including any trained $A$; its optimization only drives the gap toward its equilibrium
value $\tfrac12(s^*-a)^2$ on the aligned family (fgib\_degeneration.tex, its sweep
proposition) --- the legitimate alignment channel, not a tightness lever.
\end{remark}

\begin{remark}[Unconstrained side head (the implementation)]\label{rem:orth}
Proposition~\ref{prop:inverse} holds for any invertible $A$: the KL invariance is exact,
and the object claim needs only that the perturbation covariance
$\sigma_p^2(A^{\top}A)^{-1}$ be \emph{known} ($A$ is a model parameter). The diagonal
variational family is not closed under $A^{-1}$ (fgib\_linear.tex, Remark on family
closure), so within the diagonal parameterization of the codebase the equivalence is
exact at the distribution level and approximate as a variational identity; a
full-covariance side head would make it exact within the family. No orthogonality
constraint on $A$ is required --- the orthogonality of the method lives in the fixed
anchor frame and in the equilibrium it induces (fgib\_degeneration.tex, aligned family).
\end{remark}

\begin{remark}[Anchor-frame reading]\label{rem:frame}
With the orthonormal anchor matrix $V$ (columns $v_y$), the transform $V^{\top}$ expresses
$z$ in class coordinates: the anchors become $a\,e_y$, and the gap of
Lemma~\ref{lem:gap} decomposes into $K$ independent per-class terms. The orthogonal frame
is the geometric content of the prior; both inverse transforms ($A^{\top}$ and $V^{\top}$)
preserve it.
\end{remark}

\begin{remark}[Positioning against VIB and NIB]\label{rem:position}
VIB's rate bounds the \emph{full} mutual information $I(X;Z)$ with a marginal prior;
CEB removes the retained label information ($\mathrm{ReX} = R - I(Y;Z)$ at the optima),
and FGIB inherits the residual-level object while replacing the learnable backward
encoder with a fixed geometric prior (Proposition~\ref{prop:ordering}).
NIB's non-parametric bound $\widehat I \ge I(X;M)$ is $y$-unconditioned and tight only
when the components share a covariance and form well-separated clusters (Kolchinsky et
al.\ 2017, Eq.~10 and the discussion following Eq.~11); when the noise is small relative
to the cluster spread it saturates at $\ln N$ and its gradient vanishes (their
Section~IVA) --- loose exactly in the regime where the parametric FGIB bound is
well-behaved. The two bounds are complementary: NIB certifies the full compression
non-parametrically when the geometry is favorable; FGIB certifies the residual
parametrically with a geometry of its own choosing.
\end{remark}

\begin{remark}[Degenerate anchors]\label{rem:a0}
At $a = 0$ the anchors collapse to $N(0, I)$ for all classes and the gap loses its class
structure: $\Delta = \mathbb{E}_y[\operatorname{KL}(q(z|y)\|N(0,I))]$ --- a VIB-style
marginal gap, and the object-tightness argument of
Proposition~\ref{prop:object}(b) still holds, but the compressed endpoint becomes the
trivial collapse (fgib\_degeneration.tex, Remark on the endpoint).
\end{remark}
